\section{Chương~\ref{sec:sleep}. Giấc ngủ}

Các chế độ ngủ, việc phát lại và sự thu nhỏ khớp thần kinh đã được đo bằng ghi nơ-ron, vi thẩm tách và hiển vi điện tử; lịch ngủ và dao động chậm được mô tả bằng các mô hình của sách, còn mục đích của giấc ngủ phần lớn là giả thuyết.

{\footnotesize
\setlength{\tabcolsep}{3pt}\renewcommand{\arraystretch}{1.15}
\begin{longtable}{>{\raggedright}p{0.5cm}>{\raggedright}p{4.0cm}>{\raggedright}p{5.6cm}>{\raggedright}p{2.3cm}>{\raggedright\arraybackslash}p{1.7cm}}
\toprule
TT & Khẳng định & Thí nghiệm và điều đã đo & Nguồn & Trạng thái \\
\midrule\endfirsthead
\toprule
TT & Khẳng định & Thí nghiệm và điều đã đo & Nguồn & Trạng thái \\
\midrule\endhead
1 & Khi ngủ, nơ-ron vỏ não không im lặng; cái thay đổi là chế độ làm việc & Ghi dài hạn nơ-ron vỏ não ở chuột cống: trong giấc ngủ sóng chậm, tần số vẫn khác không, các giai đoạn làm việc xen kẽ với các giai đoạn im lặng; sau thời gian thức dài, tần số cao hơn ở mọi trạng thái, sau khi ngủ thì thấp hơn & \cite{Vyazovskiy2009} & đã đo \\
2 & \nt{ACET} cao ban ngày và trong giấc ngủ REM, thấp trong giấc ngủ sóng chậm (bảng~\ref{tab:sleepmodes}) & Vi thẩm tách trong vỏ não ở mèo tự do vận động: khi thức, lượng acetylcholine giải phóng cao hơn $100$--$175\%$ so với khi ngủ sóng chậm, còn trong giấc ngủ REM thì xấp xỉ mức thức yên tĩnh & \cite{Marrosu1995} & đã đo \\
3 & \nt{NORA} và \nt{SERO} lắng xuống trong giấc ngủ sóng chậm và gần như im lặng trong giấc ngủ REM & Ghi từng nơ-ron \nt{NORA} ở chuột cống và nơ-ron \nt{SERO} ở mèo theo các giai đoạn ngủ: tần số cực đại khi thức, thấp hơn trong giấc ngủ sóng chậm và gần bằng không trong giấc ngủ REM & \cite{AstonJonesBloom1981,McGintyHarper1976} & đã đo \\
4 & Trong giấc ngủ REM, cơ bị tắt bằng ức chế qua \nt{GABA} và \nt{GLYC} & Ở chuột cống, người ta đo hoạt động của các nơ-ron \nt{ACET} của cơ nhai và đưa chất chẹn trực tiếp tới chúng: tình trạng liệt chỉ được giải trừ nếu chẹn cả hai loại thụ thể \nt{GABA} lẫn thụ thể \nt{GLYC} & \cite{BrooksPeever2012} & đã đo \\
5 & Áp lực ngủ $S$ là một tụ điện với hai ngưỡng~\eqref{eq:sleeppressure}, $\tau_r=18.2$~giờ, $\tau_d=4.2$~giờ & Mô hình hai quá trình; các hằng số thời gian thu được từ cách công suất sóng chậm của EEG tăng khi thức và giảm khi ngủ, còn dạng của các ngưỡng từ thời điểm tự thức dậy & \cite{Borbely1982,Daan1984} & lý thuyết \\
6 & Cỗ máy tự đi tới $16.1$~giờ thức và $8.0$~giờ ngủ (hình~\ref{fig:twoprocess}) & Tính theo~\eqref{eq:sleeppressure} với ngưỡng dao động, tập lệnh \texttt{models/sleep\_models.py}: $16.1$~giờ thức và $7.9$--$8.0$~giờ ngủ & --- & mô hình \\
7 & Sau $40$~giờ không ngủ $S=0.90$, nhưng giấc ngủ chỉ kéo dài $8.8$~giờ: món nợ được trả bằng độ sâu, không bằng độ dài & Mô hình: \texttt{models/sleep\_models.py} cho $S=0.90$ và $8.8$~giờ. Thí nghiệm: sau $40.5$~giờ không ngủ, ở tám người, trong đêm ngủ bù đầu tiên công suất sóng chậm của EEG cao hơn đáng kể so với bình thường & \cite{Borbely1981} & mô hình \\
8 & Về mặt vật lý, $S$ là adenosine & Vi thẩm tách ở mèo: adenosine ngoại bào ở một trong các vùng điều khiển sự thức tăng lên trong thời gian thức, tăng tiếp khi bị thiếu ngủ và giảm khi ngủ. Việc $S$ chỉ là adenosine chưa được chứng minh & \cite{PorkkaHeiskanen1997,Dunwiddie2001} & giả thuyết \\
9 & Trong giấc ngủ sóng chậm, vỏ não dao động: làm việc vài trăm ms rồi im lặng, chưa tới một lần mỗi giây ($<1$~Hz, khi gây mê thường $0.3$--$0.4$~Hz) & Ghi nội bào ở mèo được gây mê: khử cực $0.8$--$1.5$~s và tăng phân cực kéo dài, nhịp $<1$~Hz, thường gặp nhất $0.3$--$0.4$~Hz; cùng sự xen kẽ làm việc và im lặng này thấy được trong giấc ngủ tự nhiên (dòng~1) & \cite{Steriade1993,Vyazovskiy2009} & đã đo \\
10 & Dao động là một nhóm có phản hồi và mệt mỏi~\eqref{eq:updown}; với $b=5$ chu kỳ $1.1$~s (giá trị mô hình) và làm việc khoảng $35\%$ thời gian, với $b=1$ làm việc đều đặn (hình~\ref{fig:updown}) & Tính toán, tập lệnh \texttt{models/sleep\_models.py}: chu kỳ $1.11$~s, tỉ lệ thời gian làm việc $0.35$ (trung bình trên một số nguyên chu kỳ); với $b=1$ tỉ lệ làm việc là $1.0$. Chu kỳ của mô hình ngắn hơn chu kỳ đo được (dòng~9) & --- & mô hình \\
11 & \nt{ACET} tắt dao động và chuyển vỏ não sang làm việc đều đặn & Kích thích nhân nơ-ron \nt{ACET} ở mèo chặn nhịp chậm ở $79\%$ nơ-ron vỏ não và chuyển chúng sang làm việc đều đặn, chủ yếu do các đợt tăng phân cực kéo dài biến mất. Acetylcholine ức chế các dòng kali chậm gây ra mệt mỏi & \cite{SteriadeAmzica1993,McCormick1992} & đã đo \\
12 & Các đợt bùng nhịp $7$--$15$~Hz do vòng \nt{GLUT}--\nt{GABA} giữa vỏ não và nút chuyển tiếp sinh ra & Trong lát nút chuyển tiếp thị giác của chồn sương, các đợt bùng được tạo ra bởi các chùm xung đáp trả của tế bào chuyển tiếp trước sự ức chế từ một nhân \nt{GABA} lân cận, mà chính chúng kích thích & \cite{vonKrosigk1993,SteriadeMcCormickSejnowski1993} & đã đo \\
13 & Việc phát lại ngày đã qua diễn ra có tua nhanh: ở hồi hải mã nhanh hơn khoảng $20$ lần, ở vỏ não $6$--$7$ lần & Trong giấc ngủ sóng chậm, các chuỗi lặp lại của nơ-ron hồi hải mã bị nén theo thời gian. Sau khi chạy trên đường chạy, các chuỗi nơ-ron vị trí được lặp lại theo cùng thứ tự trong các đợt ngắn $\sim100$~ms, nén khoảng $20$ lần; ở vỏ não, mức nén là $6$--$7$ lần & \cite{Nadasdy1999,LeeWilson2002,Euston2007} & đã đo \\
14 & Tăng cường sự phối hợp giữa phát lại và vỏ não cải thiện trí nhớ (quan hệ nhân quả) & Ở chuột cống sau khi học, kích thích gắn với các lần phát lại làm tăng liên kết của chúng với sóng chậm và các đợt bùng nhịp ở vỏ não: hôm sau trí nhớ ở mức cao, còn ở nhóm đối chứng thì ở mức ngẫu nhiên & \cite{Maingret2016} & đã đo \\
15 & Mục đích của phát lại là học trộn lẫn cho bộ nhớ chậm & Lý thuyết hai hệ thống học và các tính toán trên mạng nhân tạo: học tuần tự làm hỏng cái cũ, học trộn lẫn thì không. Chưa có thí nghiệm trực tiếp trên não cho thấy phát lại cần chính là cho việc này & \cite{McClelland1995} & giả thuyết \\
16 & Sau khi ngủ, khớp thần kinh thu nhỏ tỉ lệ với kích thước, các khớp lớn hầu như không đổi & Hiển vi điện tử ba chiều $6920$ khớp thần kinh vỏ não chuột nhắt: diện tích tiếp xúc sau khi ngủ nhỏ hơn $\sim18\%$ so với sau khi thức, mức giảm tỉ lệ với kích thước, các khớp lớn và bền không thay đổi & \cite{deVivo2017} & đã đo \\
17 & Nhiệm vụ chính của giấc ngủ là chuẩn hóa lại trọng số & Giả thuyết cân bằng nội môi khớp thần kinh; các dòng 1 và 16 phù hợp với nó, nhưng không chứng minh đó là chức năng chính & \cite{TononiCirelli2014} & giả thuyết \\
18 & Giấc ngủ REM là việc thử nghiệm mô hình thế giới trên bàn thử & Chế độ (bảng~\ref{tab:sleepmodes}) đã được đo; việc mạng chạy thử mô hình thế giới là một giả định lý thuyết, chưa có thí nghiệm & \cite{Hobson2009} & giả thuyết \\
19 & Rửa sạch não khi ngủ: câu hỏi còn mở & Một thí nghiệm: khi ngủ, khoảng gian bào rộng hơn $60\%$ và việc dọn dẹp nhanh hơn. Một thí nghiệm khác, với cách đo khác: khi ngủ và khi gây mê, việc dọn dẹp chậm hơn rõ rệt. Các kết quả mâu thuẫn nhau & \cite{Xie2013,Miao2024} & giả thuyết \\
20 & Giấc ngủ cục bộ: các vùng vỏ não của con vật đang thức thỉnh thoảng rơi vào im lặng trong chốc lát & Ở chuột cống sau thời gian thức dài, nơ-ron của một vùng vỏ não riêng lẻ im lặng trong chốc lát như trong giấc ngủ sóng chậm, kèm một sóng chậm trong EEG cục bộ; vào những lúc đó chuột mắc lỗi trong nhiệm vụ nhiều hơn & \cite{Vyazovskiy2011} & đã đo \\
\bottomrule
\end{longtable}}
